%% SCRAP: archive/session-logs/experimental-addendum
%% SOURCE: docs/working/archive/session-logs/experimental-addendum.md
%% STATUS: HISTORICAL
%% FITS: experiments/ if content warrants
%% EDITORIAL: lifted — prose rewritten to press voice. Source is a
%%   ChatGPT transcript; the scientific content on James Law and the
%%   window-scaling experiment design is the load-bearing material
%%   and is retained. Informal dialogue omitted.

\section{Experimental Addendum: James Law and Window-Scaling Validation}
\label{sec:experimental-addendum}

\subsection{Discovery Context}

During analysis of the SSM attractor results, a conserved geometric
relationship was identified in the feedback-response surface of the
adaptive runtime. The system exhibits:

\begin{itemize}
  \item Reproducible convergence to a stable attractor basin.
  \item Shape-invariant timing behaviour across workloads.
  \item A consistent geometric structure in the configuration--window--variance
    phase space.
  \item A predictable collapse when the smoothing window exceeds a critical
    bound.
\end{itemize}

These properties distinguish the StarForth adaptive runtime from chaotic
adaptive systems. The behaviour is geometric and predictable, not
stochastic.

\subsection{James Law of Computational Dynamics}

The empirically observed scaling relation is stated as follows.

\begin{quote}
\textbf{James Law.} In an adaptive computational system governed by
multiple interacting feedback loops, the effective stability-smoothing
factor $\Lambda$ scales inversely with the number of active degrees of
freedom $\mathrm{DoF}$ plus one, and directly with the smoothing window
size $W$:
\[
  \Lambda = \frac{W}{\mathrm{DoF} + 1}
\]
This relation describes a conserved geometric structure in the system's
feedback-response surface and predicts the onset of stability,
metastability, and collapse phases as window size varies.
\end{quote}

%% PATENT: This section is adjacent to patent claims. Do not expand or
%% reformulate without explicit instruction.

The law is currently a working hypothesis. Prior to incorporation into
any patent application, it is to be validated empirically via the
window-scaling experiment described below.

\subsection{Scientific Significance}

If validated, the James Law occupies the same class of discovery as
Little's Law in queueing theory, Amdahl's Law in parallelism, and
Zipf's Law in linguistics — an empirically observed, falsifiable,
reproducible scaling invariant in a new domain (computational dynamics).

\subsection{Window-Scaling Validation Experiment}

The validation experiment tests whether the quantity

\[
  K = \frac{\Lambda(\mathrm{DoF}) \cdot (\mathrm{DoF} + 1)}{W}
\]

remains approximately constant and near unity across multiple degrees of
freedom and window sizes.

\subsubsection{Parameters}

\begin{itemize}
  \item \textbf{Degrees of freedom:} $\mathrm{DoF} \in \{0, 1, 2, 3, 4, 5, 6, 7\}$.
  \item \textbf{Window sizes:} $W \in \{512, 1024, 1536, 2048, 3072, 4096,
    6144, 8192, 16384, 32769, 52153, 65536\}$ (12 levels, spanning
    subcritical through catastrophic-collapse territory).
  \item \textbf{Replicates:} 30 per condition.
  \item \textbf{Total runs:} $8 \times 12 \times 30 = 2{,}880$.
  \item \textbf{Workload:} composite ``omni'' workload — all L8 initialisation
    scripts concatenated and run sequentially
    (STABLE $\to$ TEMPORAL $\to$ VOLATILE $\to$ TRANSITION $\to$ DIVERSE),
    maximally stressing the feedback loops.
\end{itemize}

\subsubsection{Run Matrix}

Conditions are generated by a full factorial expansion and shuffled once
before execution (fixed random seed for reproducibility) to eliminate
temporal drift, thermal bias, and cache-warming artefacts.

\subsubsection{Analysis Criteria}

\begin{enumerate}
  \item Compute $K(W, \mathrm{DoF})$ for each condition.
  \item Measure mean, standard deviation, and maximum deviation of $K$
    from~1.0.
  \item Identify the critical window size $W_\mathrm{crit}(\mathrm{DoF})$
    at which the CV spikes, the L8 engine collapses to a trivial mode,
    or the attractor structure disappears.
\end{enumerate}

\subsubsection{Acceptance Criteria}

The law is supported if:
\begin{itemize}
  \item $K \in [0.97, 1.03]$ across all DoF levels at $W = 4096$.
  \item The same clustering holds at $W = 2048$ and $W = 8192$.
  \item Instability (elevated CV or mode collapse) correlates with
    deviation from $K \approx 1$ at extreme window sizes.
\end{itemize}

\subsection{Presentation Guidance}

When discussing the James Law in publications, the analogy to physical
systems should be used as interpretive framing only, not as a claim.
Appropriate language:

\begin{quote}
``While the runtime does not implement gravitational physics, the emergent
topology resembles a system with a dominant attractor. This analogy is
useful as a descriptive shorthand, but the formal behaviour is defined
strictly by the equations presented herein.''
\end{quote}

The experimental evidence — DoE data, convergence plots, variance
shrinkage, attractor identification — constitutes the scientific
claim. The metaphors motivate the mathematics; they do not replace it.
